\begin{problem}{\bf Linear (and affine) encryption schemes are insecure} (10 points)\\
Let $n \geq 2$ be an integer. Let $A$ be an invertible
$n \times n$ binary matrix, and let $b$ be a $1 \times n$ binary
vector. In the Hill symmetric-key encryption scheme, the secret
key is a pair $(A,b)$. Plaintext messages $m$ are represented as
$1 \times n$ binary vectors. The encryption function is
$E(m) = mA + b$.\\
\emph{Note: All arithmetic is performed modulo 2. You may assume
that linear algebra ``works'' modulo 2.}\\
Here is a toy example with $n=5$:
\[ A = \left[
       \begin{array}{ccccc}
       1&0&0&0&1\\
       0&1&0&1&1\\
       0&0&1&1&0\\
       0&1&0&1&0\\
       0&0&1&0&1\\
       \end{array} \right] ,
   \;\;\;\;
   b = [1,1,0,1,1],
   \;\;\;\;
   m = [1,0,1,0,1],
   \;\;\; \mbox{and } \;\;\; E(m) = [0,1,0,0,1].
\]
\begin{enumerate}[label=(\alph*)]
\item Give a decryption formula for the Hill cipher (with general
      parameter $n$).
\item Show how an adversary can efficiently determine the secret key
      $(A,b)$ using a \emph{chosen-plaintext} attack.
\item In a \emph{chosen-ciphertext} attack, the adversary is given a
      challenge ciphertext $c$. She can obtain (by asking Alice or Bob,
      who know the secret key) the plaintext of \emph{any} ciphertext of
      her choice \emph{except for $c$ itself}. Her task is to decrypt $c$.\\
      Show how the adversary can decrypt $c$ by obtaining the decryptions
      of at most 3 ciphertexts (each different from $c$) from Alice or Bob.
\end{enumerate}
\end{problem}

\begin{loigiai}
\end{loigiai}
